\documentclass[11pt]{article}
\usepackage[margin=2.5cm]{geometry}
\usepackage{xcolor,enumitem,amsmath}
\newcommand{\E}{\mathrm{E}}
\setlength{\parindent}{0pt}
\setlength{\parskip}{0.5em}
\setlength{\emergencystretch}{2em}
% Set to \draftfalse ONLY after the student has confirmed (and edited) the text below.
\newif\ifdraft
\drafttrue

\newcommand{\draftbanner}{\ifdraft\noindent\fcolorbox{red}{red!8}{\parbox{\dimexpr\linewidth-2\fboxsep-2\fboxrule}{\color{red!70!black}\small
\textbf{DRAFT -- REQUIRES STUDENT CONFIRMATION.} This text describes the proposed division of work. The student must
check every item, change it so that it describes what they actually did (including, where relevant, which parts were
prepared with the help of an AI assistant), fill in every red placeholder, and then switch off this banner
(\texttt{workdone/status.tex}) before submission.}}\par\bigskip\fi}

% Placeholder for text only the student can write; package.sh refuses final packaging while one remains.
\newcommand{\fillin}[1]{\textcolor{red}{[STUDENT TO FILL IN: #1]}}

\begin{document}
\draftbanner
I worked on the following parts in this project:

\textbf{Mert Seyhan (2023400213) -- CMPE 300 Project 1, variant key 4}

\begin{enumerate}[leftmargin=1.5em]
\item \textbf{Reading the algorithm and the counting conventions.} I fixed how the marked statements are counted:
a comparison is counted each time it is evaluated (line 8 only when line 3 is false), and a \texttt{for} loop is
expanded into its \texttt{while} form, so the loop on line 9 makes $s+1$ increments and the loop on line 15 makes
$n-i$ increments. I computed the variant key $4$ ($p_a=p_b=2/5$, $p_c=1/5$, $d=2$) and wrote Section 2 of the report.
\item \textbf{Theoretical analysis of group (1).} I derived $\tau_1=2n-n_a$ and from it $B_1(n)=n$ (input
$\texttt{a}^n$), $W_1(n)=2n$ (input $\texttt{b}^n$) and $A_1(n)=8n/5$, with the per-iteration proofs and the
$\Theta(n)$ bounds (Section 4).
\item \textbf{Theoretical analysis of group (2).} I proved $B_2(n)=0$ and $W_2(n)=n(n+1)/2$ with worst input
$\texttt{c}^n$, using the pair-counting bound $P\le\sum_{arr[j]=\texttt{a}}(n-1-j)$ to show that no input with
\texttt{a}'s and \texttt{b}'s can beat $\texttt{c}^n$. I derived $\E[S_i]=2(1-(4/5)^i)$ from the recurrence
$e_{i+1}=\frac45e_i+\frac25$ (backward substitution), its sum, and $A_2(n)=\frac{n^2}{10}+\frac{13n}{10}-4+4(\frac45)^n$
with the $\Theta(n^2)$ constants (Section 5). The $\E[S_i]$ result is also used for $A_3$ and $A_4$.
\item \textbf{Python implementation (\texttt{2023400213.py}).} I wrote the timing version of the algorithm without
counters, the counted version with the four counters, the random input generator, the best/worst input generators of
group (4), and the command-line interface (normal mode and \texttt{--count} mode, argument checks, exact output
format).
\item \textbf{Timing infrastructure and Table 2.} I wrote the timing function (\texttt{perf\_counter}, garbage
collector off, repetition until 20\,ms, minimum of 3 rounds), ran the 45 measurements with 10 random runs per average,
recorded the machine information, and filled Table 2. I noticed the noisy first run and documented why the 3-round
rule was introduced (Section 10).
\item \textbf{Choice of the basic operation.} I wrote Section 9.2: the argument that group (4) counts the bodies of
all innermost loops, the measured time per operation of each group, and Table A1 with the extreme inputs of every
group at $n=400$.
\item \textbf{Tests.} I wrote the tests for the variant, the command-line interface (subprocess tests of both modes and
of wrong arguments) and the timing code.
\item \textbf{Cross-checking.} \fillin{which parts of Enes's work (e.g.\ the derivations of $W_3(n)$ and $W_4(n)$)
you checked, and how}. I read the final report.
\end{enumerate}
\end{document}
